% Consolidated guide-ch15 from f8f4709b0ffac9da4c1fecbf2ce873002f6da6cc; source fragments preserved.

% === BEGIN guide-ch15-examples.tex blob cbd624456880fb5c76d5bdd20730a8a61fb97c64 ===
% 第13批：《大学物理学学习指导（第5版）》第15章“量子物理基础”
% 例15.1—15.8；印刷220—223页／PDF物理228—231页。
% 全部原题数字逐页核对；特别保留15.2书内康普顿波长指数错误。
\subsection{例题：光电效应、康普顿散射、物质波与势阱}

\sourceexample{15.1}{220}{228}
\problemtext 同一种金属发生光电效应，
入射单色光波长由$\lambda_1=400$ nm缩短至
$\lambda_2=360$ nm。
求遏止电压的改变量及变化方向。
\workings
对于同一种光电阴极，逸出功$W$恒定，
由$eU_a=hc/\lambda-W$，两次作差消去$W$，
\[
 \Delta U_a=\frac{hc}{e}\left(\frac1{\lambda_2}
                  -\frac1{\lambda_1}\right).
\]
用$hc/e\simeq1239.842\,\mathrm{V\,nm}$计算，
\[
 \boxed{\Delta U_a=
 1239.842\left(\frac1{360}-\frac1{400}\right)
 \,\mathrm V\approx+0.3444\,\mathrm V}.
\]
所以遏止电压的\textbf{绝对值增大约0.345V}。
\crosscheck 波长缩短，光子能量增加，
因此最大光电子动能和所需遏止电压绝对值均增加。
若仪器将遏止电压记录为负偏压，带符号示数反而变得更负。

\sourceexample{15.2}{221}{229}
\problemtext 康普顿散射中，入射X射线波长
$\lambda_0=3.0\times10^{-3}$ nm，
反冲电子速率$v=0.6c$，
假设电子初始静止，求散射光子波长及散射角。
\workings
电子反冲的相对论动能为
\[
 K_e=(\gamma-1)m_ec^2,\qquad
 \gamma=(1-0.6^2)^{-1/2}=1.25,
\]
即$K_e=0.25m_ec^2\simeq127.75$ keV。
初始光子能量
$E_0=hc/\lambda_0\simeq413.28$ keV。
由光子与电子的总能量守恒，
散射后光子能量
$E'=E_0-K_e\simeq285.53$ keV，因而
\[
 \boxed{\lambda'=\frac{hc}{E'}
 \simeq4.342\times10^{-3}\,\mathrm{nm}}.
\]
进一步利用康普顿位移式
\[
 \lambda'-\lambda_0
 =\lambda_C(1-\cos\varphi),\quad
 \lambda_C=\frac{h}{m_ec}\simeq2.42631\times10^{-3}\,\mathrm{nm},
\]
得到
\[
 \boxed{\cos\varphi
 =1-\frac{\lambda'-\lambda_0}{\lambda_C}},
 \qquad
 \boxed{\varphi\simeq63.5^\circ}.
\]
\crosscheck \textbf{原书数值存在明显自洽性问题：}
源书印刷页221将散射波长写成
$4.3\times10^{-4}$ nm，竟短于入射的
$3.0\times10^{-3}$ nm，违背自由静止电子散射时
$\lambda'\ge\lambda_0$。
依据源题数据守恒计算应为
$4.34\times10^{-3}$ nm，
相应散射角约$63.5^\circ$；
书上写约$62^\circ24'$，
可能是中间过程仅保留两位有效数字所致。
本稿不默默照抄矛盾的指数。

\sourceexample{15.3}{221--222}{229--230}
\problemtext 氢原子被外来可见单色光激发后，
在它的巴耳末发射系中仅见三条不同谱线。
求这三条光谱的波长和外来激发光的频率
（按书中的玻尔/里德伯主能级模型）。
\workings
巴耳末发射线都终止于$n=2$。
若出现且仅出现由较高主量子数发向$n=2$的三条，
则最高激发主量子数为$n=5$：
$n=5,4,3\to2$。
里德伯公式给
\[
 \frac1{\lambda_{n\to2}}=
 R_H\left(\frac14-\frac1{n^2}\right),
 \qquad R_H\simeq1.097\times10^7\,\mathrm{m^{-1}}.
\]
于是
\[
 \boxed{\lambda_{5\to2}\approx434\,\mathrm{nm}},\quad
 \boxed{\lambda_{4\to2}\approx486\,\mathrm{nm}},\quad
 \boxed{\lambda_{3\to2}\approx656\,\mathrm{nm}}.
\]
源题明确指出\emph{外来激发光是可见光}，
故按所给模型，原子须由已有的$n=2$态吸收
$\Delta E=E_5-E_2$进入$n=5$，
而不能由基态$n=1$直接吸收同一可见光：
\[
 \boxed{\nu_{\rm pump}
 =cR_H\left(\frac14-\frac1{25}\right)
 \approx6.91\times10^{14}\,\mathrm{Hz}}.
\]
\crosscheck 不要混同“激发到$n=5$”与
“从基态$n=1$激发到$n=5$”。
用可见光完成激发必须有合适的初始激发态；
更精确的原子跃迁还受角动量选择定则约束，
这里仅沿用教材能级题模型。

\sourceexample{15.4}{222}{230}
\problemtext 判断“波由粒子构成”与“粒子由波构成”
两个关于德布罗意波的说法是否正确。
\workings
两种说法若理解为经典粒子集合与经典连续波的
简单机械拼装，均不成立。
一个一个电子依次通过双缝时，
每次探测为局域事件，累积的概率分布却出现干涉条纹。
适宜的量子理论表述是由量子态或波函数
确定测量结果的概率幅，
\[
 \boxed{P(x)\,dx=|\psi(x)|^2dx}
\]
（按一维、已归一化态举例）。
\crosscheck 概率幅不能当作“由经典波组成的小球”；
单个检测点也不意味着电子在测量前必须沿唯一经典轨道传播。
本题为概念题，不伪造不存在的定量数值。

\sourceexample{15.5}{222}{230}
\problemtext 比较书中两幅一维粒子波包图：
(a) 多个周期、空间延展较宽，
(b) 在$x$方向集中成一个狭窄脉冲。
分别判断位置与动量的确定程度。
\workings
由傅里叶波包的频谱--宽度关系，
空间展开得越宽，波数扩展越窄；
由德布罗意关系$p=\hbar k$，动量分布也越窄。
因此(a)的空间展宽大于(b)，而动量分布展宽较小：
\[
 \boxed{\Delta x_a>\Delta x_b},\qquad
 \boxed{\Delta p_{x,a}<\Delta p_{x,b}}.
\]
更严格地，标准差定义下
\[
 \boxed{\Delta x\,\Delta p_x\ge\frac{\hbar}{2}}.
\]
\crosscheck 同一瞬间图像给出的是波包\emph{包络的空间尺度}，
不是沿$x$传播的粒子“经典轨迹长度”。
两幅图不提供定量高斯宽度，
不应该凭线条数量捏造$\Delta x$的数值。

\sourceexample{15.6}{223}{231}
\problemtext 电子动能$K=1.0$ keV，局域化于
$\Delta x=10^{-10}$ m范围内。
依据\emph{书中指定的估算关系}
$\Delta x\Delta p_x\gtrsim h$，
求动量不确定度占动量的百分比。
\workings
电子在1keV下可采用非相对论动能近似：
\[
 p=\sqrt{2m_eK}
 =\sqrt{2m_e(1000\,\mathrm{eV})}
 \simeq1.71\times10^{-23}\,\mathrm{kg\,m/s}.
\]
按原书指定的$\Delta x\Delta p_x\gtrsim h$，
\[
 \Delta p_x\gtrsim\frac h{\Delta x}
 \simeq6.626\times10^{-24}\,\mathrm{kg\,m/s},
\]
所以
\[
 \boxed{\frac{\Delta p_x}{p}\gtrsim0.388
 \approx38.8\%}.
\]
\crosscheck \textbf{这不是海森堡标准差的不等式常数。}
严格标准形式为
$\Delta x\Delta p_x\ge\hbar/2=h/(4\pi)$，
若$\Delta x$按严格标准差定义，
对应下界约$3.09\%$。
原书估算的38.8%应与其自定的\emph{宽度量级约定}绑定，
不能宣称是唯一的严格最小百分比。

\sourceexample{15.7}{223}{231}
\problemtext 一维无限深势阱的墙面为$x=0$和$x=L$。
求在$0<x<L/3$区域内找到粒子的概率，
分别对$n=1,2,3$定态以及经典匀速分布模型作比较。
\workings
势阱归一化定态
$\psi_n(x)=\sqrt{2/L}\sin(n\pi x/L)$。
积分概率密度：
\[
 \begin{aligned}
 P_n&=\int_0^{L/3}|\psi_n(x)|^2\,dx\\
 &=\frac{2}{L}\int_0^{L/3}
       \sin^2\frac{n\pi x}{L}\,dx
 =\boxed{\frac13-\frac{\sin(2n\pi/3)}{2n\pi}}.
 \end{aligned}
\]
代入得
\[
 \boxed{P_1=\frac13-\frac{\sqrt3}{4\pi}\approx0.196},
 \qquad
 \boxed{P_2=\frac13+\frac{\sqrt3}{8\pi}\approx0.402},
 \qquad
 \boxed{P_3=\frac13}.
\]
经典匀速运动且位置分布均匀时，
$\boxed{P_{\rm classical}=1/3}$。
\crosscheck 量子定态位置密度非均匀；
虽然当$n=3$时恰给$1/3$，
不能据此认定所有量子定态都等于经典均匀分布。

\sourceexample{15.8}{223}{231}
\problemtext 求一个$d$支壳层所能容纳的最大电子数，
列出相关轨道磁量子数$m_l$与自旋磁量子数$m_s$。
\workings
$d$亚层的轨道角量子数为$l=2$，
因而
\[
 \boxed{m_l=-2,-1,0,+1,+2},\qquad
 \boxed{m_s=-\tfrac12,\,+\tfrac12}.
\]
共有$2l+1=5$个轨道方向，每方向有2种自旋态。
泡利不相容原理使每个不同的$(m_l,m_s)$
组合最多被一电子占据，于是
\[
 \boxed{N_d=2(2l+1)=2(5)=10}.
\]
\crosscheck 这里的10是\emph{一个$d$亚层}最大容纳数，
而不是一个完整$n$壳层的总态数$2n^2$。
% === END guide-ch15-examples.tex ===
